\documentclass[11pt,a4paper]{article}
\usepackage[margin=2cm]{geometry}
\usepackage{graphicx}
\usepackage{booktabs}
\usepackage{amsmath}
\usepackage{siunitx}
\usepackage[hidelinks]{hyperref}
\usepackage{caption}
\usepackage{float}
\captionsetup{font=small}
\setlength{\parskip}{4pt}

\title{A conditional diffusion emulator of Veros ACC equilibrium states:\\first end-to-end results}
\author{}
\date{18 September 2026}

\begin{document}
\maketitle

\section*{Summary}
We trained a conditional denoising diffusion model to generate the post-spin-up temperature and salinity
state of the Veros ACC channel configuration directly from the two coefficients of its eddy-kinetic-energy
closure, $c_k$ and $c_\epsilon$. On ten held-out parameter pairs the eight-sample ensemble mean has an RMSE
of \SI{0.030}{K} against the true 20-year mean state, with a domain-mean bias below \SI{1}{mK}, and the
generated response over the full $10\times10$ parameter grid reproduces the true one. The choice of
normalisation decides whether conditioning works at all: DINO-Fusion's per-level scaling makes the model
collapse to the mean state (\SI{0.154}{K}), per-cell anomaly scaling does not. The emulator beats the nearest
training run (\SI{0.037}{K}) but not yet the average of the four neighbouring runs on the grid
(\SI{0.020}{K}). The full chain costs under one A100 hour.

\section{Data}
One hundred Veros ACC runs on a complete $10\times10$ log-spaced grid, $c_k\in[0.0125,\,0.8]$ and
$c_\epsilon\in[0.0875,\,5.6]$, 100 years each with 30-day output on a $30\times42\times15$ interior grid
(zonal, meridional, vertical). We keep temperature and salinity for the last 20 years of every run
(241 snapshots per run). Three properties of the ensemble shape the problem:
\begin{itemize}\itemsep1pt
  \item Salinity is exactly \SI{35}{psu} in every water cell of every run; it is carried through the
        pipeline as a second field but contains no information.
  \item The end states are nearly deterministic: after removing the residual drift, the grid-point
        temporal standard deviation is \SI{0.002}{K}, against \SI{0.05}{K} between neighbouring runs on
        the grid and up to \SI{1.5}{K} between extremes. The runs still drift by 0.03 to \SI{0.12}{K}
        in domain-mean temperature over the last 20 years, so the target is the state at year 100
        rather than a true equilibrium.
  \item The two parameters act almost as one: $\log c_k - 1.05\log c_\epsilon$ explains 99\,\% of the
        variance of the domain-mean temperature, and the low-$c_k/c_\epsilon$ part of the grid saturates
        to one state. The response lives between 600 and \SI{800}{m} and near the bottom, not at the
        restored surface.
\end{itemize}
Ten interior grid points were held out at random before any statistic was computed; the remaining 90 runs
(21\,690 snapshots) form the training set.

\section{Method}
The pipeline is a fork of the DINO-Fusion diffusion code adapted to a parameter ensemble.
Temperature and salinity are stacked into 30 channels, normalised, set to zero on the fixed land mask, and
zero-padded from $42\times30$ to $48\times32$ so the UNet can halve the grid four times. Two normalisations
were trained with otherwise identical settings:
\begin{align*}
  \text{per level (DINO):}\quad & x' = (x-\mu_z)/(3\sigma_z), \\
  \text{per cell (anomaly):}\quad & x' = (x-\mu_{zyx})/(3\max(\sigma_{zyx},\,0.05)),
\end{align*}
with statistics from the training runs only. The floor keeps the constant salinity at exactly zero.
The denoiser is a four-level UNet2DModel (64, 64, 128, 128 channels, 6.9\,M parameters). The standardised
$(\log c_k,\log c_\epsilon)$ pair passes through a small MLP whose output is added to the timestep
embedding. Training is the standard DDPM $\epsilon$-prediction objective with 1000 steps, a squared-cosine
schedule, sample clipping, AdamW at $10^{-4}$ with cosine decay, batch 32, 20\,000 steps and an EMA of the
weights. Sampling runs the 1000-step DDPM loop with the EMA weights and re-imposes zeros on land and padding
after every step.

\section{Evaluation}
For each held-out run we draw eight samples at its $(c_k,c_\epsilon)$ and compare, on water cells, with the
true time-mean state of the last 20 years. Baselines: the mean training state, the nearest training run in
standardised log-parameter space, and the average of the axis-neighbour training runs on the grid.
We also generate four samples for every one of the 100 grid points to compare the map of domain-mean
temperature with the truth.

\begin{table}[H]
\centering\small
\begin{tabular}{lrr}
\toprule
 & anomaly & per level \\
\midrule
diffusion, ensemble mean of 8: RMSE (K) & $0.030 \pm 0.009$ & $0.154 \pm 0.018$ \\
diffusion, single sample: RMSE (K) & 0.072 & 0.188 \\
neighbour average: RMSE (K) & $0.020 \pm 0.011$ & $0.020 \pm 0.011$ \\
nearest training run: RMSE (K) & $0.037 \pm 0.027$ & $0.037 \pm 0.027$ \\
mean training state: RMSE (K) & $0.155 \pm 0.030$ & $0.155 \pm 0.030$ \\
diffusion: domain-mean bias (K) & $-0.0002$ & $-0.028$ \\
ensemble spread / true within-window spread (K) & 0.049 / 0.031 & 0.105 / 0.031 \\
interfaces with $T$ decreasing upward, generated / truth & 8.4\,\% / 7.7\,\% & 9.5\,\% / 7.7\,\% \\
max salinity error (psu) & 0.028 & 0.010 \\
domain-mean $T$ map, RMSE on hold-out / all runs (K) & 0.012 / 0.045 & 0.030 / 0.049 \\
\bottomrule
\end{tabular}
\caption{Hold-out metrics, mean $\pm$ standard deviation over the ten held-out runs, water cells only.
Baselines are identical in both columns because they do not depend on the model.}
\end{table}


\section{Findings}
\begin{enumerate}\itemsep1pt
  \item \textbf{Normalisation decides whether conditioning works.} The run-to-run signal is about
        2\,\% of the spatial structure in the upper ocean, so under per-level scaling the conditional
        model sees almost no signal and reproduces the training mean; its parameter map is flat and noisy.
        Per-cell anomaly scaling makes the signal order one everywhere and the model follows the parameters,
        including the warm high-$c_k$, low-$c_\epsilon$ corner (Figure~\ref{fig:grid}).
  \item \textbf{Better than the nearest run, not yet better than grid interpolation.} The ensemble mean
        reaches \SI{0.030}{K} against \SI{0.020}{K} for the neighbour average. Single samples are
        1.6 times more dispersed than the true within-window variability.
  \item \textbf{The error sits below \SI{1200}{m}} (Figure~\ref{fig:profile}), as a warm bias of one to two tenths of a kelvin in the
        deepest layers, where the anomaly distribution has the heaviest tails and about 2\,\% of training
        values fall outside the clipping range. The cold bias at the extreme corner run is the regime change
        at the edge of the training set.
  \item \textbf{Plumbing checks pass.} Salinity returns within \SI{0.03}{psu} of 35 without any constraint,
        land is exact, and the fraction of temperature inversions matches the truth. Figure~\ref{fig:samples} shows
        random samples next to the true state at three depths.
\end{enumerate}

\begin{figure}[H]
\centering
\includegraphics[width=0.86\textwidth]{../full_anomaly/grid_domain_mean.png}\\[4pt]
\includegraphics[width=0.86\textwidth]{../full_3std/grid_domain_mean.png}
\caption{Domain-mean temperature over the parameter grid: truth, generated, and difference.
Top: anomaly normalisation. Bottom: per-level normalisation. White dots mark training runs, red crosses
the held-out runs.}
\label{fig:grid}
\end{figure}

\begin{figure}[H]
\centering
\includegraphics[width=0.28\textwidth]{../full_anomaly/rmse_profile.png}
\caption{Anomaly normalisation: RMSE against the truth by depth, mean over the held-out runs, for the
diffusion ensemble mean and the three baselines.}
\label{fig:profile}
\end{figure}

\begin{figure}[p]
\centering
\includegraphics[width=0.72\textwidth]{../full_anomaly/samples_levels.png}
\caption{Visual check for one held-out condition ($c_k=0.2$, $c_\epsilon=2.222$, anomaly model): the true
20-year mean state (left column) and three random samples at the surface, at about \SI{650}{m} and at the
bottom, for temperature (top three rows) and salinity (bottom three rows). Surface and mid-depth samples are
indistinguishable from the truth; one sample departs visibly at the bottom, where the RMSE and the spread are
largest. Salinity is \SI{35}{psu} everywhere in the truth; the salinity panels show the residual noise of the
generator on a colour scale of a few hundredths of a psu. Grey is land.}
\label{fig:samples}
\end{figure}

\section{Next steps}
Raise the anomaly scale to four standard deviations or disable clipping; draw 32 samples per condition,
since the ensemble mean is what competes with the baselines; hold out the whole $c_k=0.8$ row to test the
regime corner; try classifier-free guidance to sharpen a small conditioning signal; and add the year as a
third condition to absorb the residual drift. The longer-term target is the same pipeline on a global Veros
configuration, where active salinity, seasonal forcing and multi-century spin-up make the generative
formulation and the density constraints of DINO-Fusion relevant.

\paragraph{Reproducibility.} Branch \texttt{miniveros} of the DINO-Fusion repository, package \texttt{miniveros/}, commit
\texttt{f912e66} for the runs reported here. Extraction takes 2.5 minutes on 8 CPU cores, each training
22 minutes on one A100, generation and evaluation 3.5 minutes per run; the whole dependency chain ran on the
development queue in under one GPU hour.

\end{document}
